\Needspace{12\baselineskip}
\label{sources}
\vspace{1mm}
{\Large\bfseries\color{darkTeal}সহায়ক গ্রন্থ}\par
\vspace{0.1mm}{\color{accentOrange}\rule{\linewidth}{0.6pt}}\vspace{0.5mm}

\begin{enumerate}
\item জাতীয় শিক্ষাক্রম ও পাঠ্যপুস্তক বোর্ড। তথ্য ও যোগাযোগ প্রযুক্তি, একাদশ--দ্বাদশ ও আলিম শ্রেণি। পঞ্চম অধ্যায়।
\item মাহবুবুর রহমান। তথ্য ও যোগাযোগ প্রযুক্তি। সিসটেক পাবলিকেশন, একবিংশ সংস্করণ, সেপ্টেম্বর ২০২৫; পঞ্চম অধ্যায়।
\item অক্ষরপত্র প্রকাশনী। তথ্য ও যোগাযোগ প্রযুক্তি। পঞ্চম অধ্যায়।
\item দি রয়েল সায়েন্টিফিক পাবলিকেশন। তথ্য ও যোগাযোগ প্রযুক্তি। পঞ্চম অধ্যায়ের বহুনির্বাচনি ও সৃজনশীল প্রশ্নসংকলন।
\end{enumerate}

\vspace{0.8mm}
{\normalsize\bfseries\color{deepOrange}পরিভাষা ও নির্দেশ যাচাইয়ের সহায়ক উৎস}\par
\vspace{0.1mm}

\begin{enumerate}
\item \textbf{উচ্চমাধ্যমিক তথ্য ও যোগাযোগ প্রযুক্তি।} ২০২৬ সালের সংশোধিত সিলেবাস।
\item \textbf{Python Software Foundation।} Python 3 Documentation:
\href{https://docs.python.org/3/reference/lexical_analysis.html}{Lexical Analysis},
\href{https://docs.python.org/3/reference/expressions.html}{Expressions},
\href{https://docs.python.org/3/tutorial/controlflow.html}{Control Flow},
\href{https://docs.python.org/3/tutorial/datastructures.html}{Data Structures},
\href{https://docs.python.org/3/tutorial/modules.html}{Modules} এবং
\href{https://docs.python.org/3/library/math.html}{math Functions}।
\end{enumerate}

